\section{Transverse configurations and determinant signs}
\label{sec:transverse}

We begin the proof of Lemma~\ref{lem:rank-layer} with the configurations
whose point differences are nonsingular. Throughout this section $D=2r$
is fixed, $q$ tends to infinity through odd primes, and constants in
$O(\cdot)$ may depend on $D$. We use $b(n)=n(n+1)/2$, including $b(0)=0$.
The estimates concerning growing $D$ will be made separately in
Section~\ref{sec:tail}.

On triples with nonsingular point differences, for two or three active
incidences on a face, the sign model has value
\(2\) or \(8\) when every active pair satisfies
\(\chi(\det(X_i-X_k))=c_0\xi_{ij}\xi_{kj}\), and value \(0\) when a
match fails. For \(r>1\), we will approximate the normalized center
factor by this model in \(L^1\) on these triples, and bound that
factor uniformly there. These estimates
permit multiplication over faces sharing point matrices. The
preliminary counts fix the layer normalizations and enumerate the
candidate subspaces; restriction-sign balance supplies the final
factor \(1/2\). We then prove that the determinant signs on all
\(33\) point pairs have an asymptotically independent joint law.

\subsection{Symmetric forms over a finite field}

The discriminant sign of a symmetric form is the quadratic character of
the determinant of its induced nondegenerate form on the quotient by
its radical. A change of basis multiplies this determinant by a square,
so its character is well defined. In dimension zero the determinant is
understood to be $1$.

Rank and discriminant sign are the standard invariants of symmetric
bilinear forms over a finite field of odd order. Exact counts by these
invariants appear, for example, in
Schmidt~\cite[Section~2, Proposition~2.1]{schmidt2015}.
We give the elementary estimates needed here, with their proofs.

\begin{lemma}[Subspaces and symmetric matrices]
\label{lem:symmetric-counts}
For $0\leq k\leq n$, the number of $k$-dimensional subspaces of
$\F_q^n$ is
\begin{equation}
\label{eq:transverse-grassmann}
 \genfrac{[}{]}{0pt}{}{n}{k}_q
 =\prod_{a=0}^{k-1}\frac{q^n-q^a}{q^k-q^a}
 =q^{k(n-k)}(1+O(q^{-1})).
\end{equation}
For a uniform matrix $M\in\mathcal S_n$, with $n\geq 1$,
\[
 \P(M\text{ is nonsingular})=1-O(q^{-1}),\qquad
 \P(\chi(\det M)=\xi)=\frac12+O(q^{-1})
 \quad(\xi\in\{+1,-1\}).
\]
For $1\leq k\leq n$, the probability of rank $k$ and prescribed
discriminant sign is
\begin{equation}
\label{eq:transverse-signed-rank}
 \left(\frac12+O(q^{-1})\right)q^{-b(n-k)}.
\end{equation}
Finally, for $0\leq u\leq n$,
\begin{equation}
\label{eq:transverse-nullity}
 \P(\dim\ker M\geq u)=O(q^{-b(u)}).
\end{equation}
\end{lemma}

\begin{proof}
Count ordered linearly independent $k$-tuples in $\F_q^n$ and divide
by the number of ordered bases of $\F_q^k$ to obtain
Equation~\eqref{eq:transverse-grassmann}. Its asymptotic follows by factoring
out the displayed powers of $q$; all products have fixed length.

For the nonsingularity and sign assertions, use induction on $n$.
On the event that the leading $(n-1)$-dimensional principal minor $C$
is nonsingular, write
\[
 M=\begin{pmatrix}C&z\\z^{\mathsf T}&a\end{pmatrix},\qquad
 \det M=\det C\,(a-z^{\mathsf T}C^{-1}z).
\]
Given $C,z$, the final scalar factor is uniform on $\F_q$. It is zero
with probability $q^{-1}$ and has either nonzero character with
probability $(q-1)/(2q)$. The exceptional principal minor has probability
$O(q^{-1})$ by induction. The same argument starts at $n=1$ with the
empty principal minor.

For the rank assertion, fix a $k$-subspace $U$ with basis-column matrix
$P$. Every symmetric matrix with image $U$ is uniquely
$PRP^{\mathsf T}$ for a nonsingular symmetric $k$-by-$k$ matrix $R$.
Indeed, if $Q P$ is the identity, symmetry and the image condition give
$M=P(QMQ^{\mathsf T})P^{\mathsf T}$, and multiplication by $Q,Q^{\mathsf T}$
also proves uniqueness. The surjection $P^{\mathsf T}:\F_q^n\to\F_q^k$
identifies the nondegenerate quotient form with $R$, so the discriminant
sign is $\chi(\det R)$. Consequently the desired count is
\[
 \genfrac{[}{]}{0pt}{}{n}{k}_q\,q^{b(k)}
 \left(\frac12+O(q^{-1})\right).
\]
Divide by $q^{b(n)}$ and use
$k(n-k)+b(k)-b(n)=-b(n-k)$.

For Equation~\eqref{eq:transverse-nullity}, sum over $u$-subspaces $U$
contained in the radical. In a basis beginning with $U$, this containment
sets to zero the first $u$ rows and columns, imposing
$nu-\binom u2$ independent linear conditions. A uniform symmetric form
therefore has $U$ in its radical with probability
$q^{-nu+\binom u2}$. Multiplication by
$\genfrac{[}{]}{0pt}{}{n}{u}_q=O(q^{u(n-u)})$ gives the assertion.
\end{proof}

In particular, writing $\pi_\xi$ for the mass of the signed rank-$r$
layer in $\mathcal S_D$, we have
\begin{equation}
\label{eq:transverse-signed-mass}
 \pi_\xi=\left(\frac12+O(q^{-1})\right)q^{-b(r)},
 \qquad \pi_\xi^{-1}=2q^{b(r)}(1+O(q^{-1})).
\end{equation}
These are the normalizations in $K_\xi$.

A subspace is \emph{totally isotropic} for a symmetric bilinear form
if the form vanishes on every pair of vectors in that subspace.
The next count is the split orthogonal counterpart of the symplectic
counts in Section~\ref{sec:orbits}; see also Taylor~\cite[Exercise~11.3]{taylor1992}.

\begin{lemma}[Split forms and their maximal isotropic subspaces]
\label{lem:transverse-split}
Let $r\geq1$. A nondegenerate symmetric form on a $2r$-dimensional
space over $\F_q$
has an $r$-dimensional totally isotropic subspace if and only if its
determinant character is $c_0=\chi((-1)^r)$. Such a form is a direct
orthogonal sum of $r$ hyperbolic planes. The total number of its
$r$-dimensional totally isotropic subspaces is
\begin{equation}
\label{eq:transverse-isotropic-count}
 N_r=\prod_{a=0}^{r-1}(1+q^a)
 =2q^{r(r-1)/2}(1+O(q^{-1})).
\end{equation}
The count includes both families of maximal isotropic subspaces.
\end{lemma}

\begin{proof}
First, a nondegenerate symmetric form has an orthogonal basis. A
nonzero form has a vector of nonzero norm, since otherwise polarization
would make the form zero. Split off this nondegenerate line and continue
on its orthogonal complement.

For any $a,b\in\F_q^\times$ and $t\in\F_q$, the two sets
$\{ax^2:x\in\F_q\}$ and $\{t-by^2:y\in\F_q\}$ each have size
$(q+1)/2$ and therefore meet. Thus $ax^2+by^2=t$ is solvable. A diagonal
nondegenerate form in at least three variables consequently has a
nonzero isotropic vector: solve $ax^2+by^2=-c$ and set the third
coordinate to $1$. Given an isotropic vector $e\ne0$, choose $f$ with
$\langle e,f\rangle=1$ and replace $f$ by
$f-\langle f,f\rangle e/2$. The resulting plane has Gram matrix
$\left(\begin{smallmatrix}0&1\\1&0\end{smallmatrix}\right)$ and splits
off orthogonally. Repetition leaves an anisotropic space of dimension
at most two. In dimension two, $ax^2+by^2$ is isotropic exactly when
$\chi(-ab)=1$, or equivalently when its determinant character is
$\chi(-1)$. Each hyperbolic plane also has determinant character
$\chi(-1)$. This proves that the determinant condition is precisely
the condition that the entire $2r$-space split into hyperbolic planes.

A split form has the required isotropic subspace, spanned by one vector
from each hyperbolic pair. Conversely, if such a subspace $U$ is given,
split off a hyperbolic plane using $e\in U$. Its orthogonal complement
has an isotropic $(r-1)$-subspace: for $u\in U$, replace $u$ by
$u-\langle u,f\rangle e$. This linear map on $U$ has kernel
$\langle e\rangle$, so its image has dimension $r-1$.
Induction shows that the entire form is split.

In split coordinates $(x,y)\in\F_q^r\oplus\F_q^r$, isotropy is the
equation $x^{\mathsf T}y=0$. There are
\[
 (q^r-1)q^{r-1}+(q^r-1)
 =(q^r-1)(q^{r-1}+1)
\]
nonzero isotropic vectors. The quotient $\ell^\perp/\ell$ by an
isotropic line has dimension $2r-2$. Splitting off a hyperbolic plane
through $\ell$ identifies this quotient with its nondegenerate
orthogonal complement, whose determinant character is
$\chi((-1)^{r-1})$. The criterion just proved makes the quotient split.
Count pairs $(\ell,U)$ with
$\ell\subset U$ and $\dim U=r$ to get
\[
 N_r\frac{q^r-1}{q-1}
 =\frac{(q^r-1)(q^{r-1}+1)}{q-1}N_{r-1}.
\]
With $N_0=1$, this gives the product in
Equation~\eqref{eq:transverse-isotropic-count}. In particular $N_1=2$;
no choice of one isotropic family has been imposed in this count.
\end{proof}

The isotropic-subspace count does not yet determine the sign of a
second form on those subspaces. The following averaged estimate
supplies that sign information for any sufficiently large family.

\begin{lemma}[Concentration of restriction signs]
\label{lem:transverse-restrictions}
Let $\mathcal U$ be any nonempty family of distinct $r$-subspaces of
$\F_q^D$, with $r\geq1$, and let $Q$ be a uniform symmetric form.
For $\xi\in\{+1,-1\}$ put
\[
 f_\xi(Q)=\frac1{|\mathcal U|}\sum_{U\in\mathcal U}
 \1_{\{Q|_U\text{ nonsingular of sign }\xi\}}.
\]
Uniformly over the family $\mathcal U$,
\begin{equation}
\label{eq:transverse-restriction-error}
 \E\left|f_\xi(Q)-\frac12\right|
 =O\left(q^{-1/2}+|\mathcal U|^{-1/2}\right).
\end{equation}
The same bound holds after conditioning $Q$ to be nonsingular, or
conditioning both $Q$ and $Q-L$ to be nonsingular, uniformly over every
fixed symmetric form $L$.
\end{lemma}

\begin{proof}
Restriction onto a fixed subspace is a surjective linear map on spaces
of symmetric forms, so each summand has expectation
$1/2+O(q^{-1})$ by Lemma~\ref{lem:symmetric-counts}.
For distinct $U,V\in\mathcal U$ write $T=U\cap V$. Extend a basis of
$T$ separately to bases of $U$ and $V$, and then to an ambient basis.
In this basis, the coefficients defining the two restricted forms
overlap exactly in the symmetric block on $T$. Thus any two forms on
$U,V$ agreeing on $T$ extend simultaneously to the ambient space;
given $Q|_T$, their remaining coefficients are independent and uniform.

If $T$ has positive dimension, its restriction is nonsingular outside
an event of probability $O(q^{-1})$; for $T=0$ this holds with probability
one. Given a nonsingular restriction on $T$, the determinant on $U$
is the determinant on $T$ times that of a uniform symmetric Schur
complement of size $r-\dim T>0$. The corresponding Schur complement
on $V$ is independent. Both signs are balanced up to $O(q^{-1})$,
uniformly over the fixed common block and mixed coefficients. The
covariance of the two indicator summands is therefore $O(q^{-1})$.
The diagonal covariance is at most $1$, giving
\[
 \Var(f_\xi)=O(q^{-1}+|\mathcal U|^{-1}).
\]
Cauchy--Schwarz and the expectation estimate prove
Equation~\eqref{eq:transverse-restriction-error}.

For the last assertion, each of $Q$ and $Q-L$ is marginally uniform.
The event $\mathcal A$ that the specified forms are nonsingular has
probability $1-O(q^{-1})$, uniformly in $L$. Since the error being
averaged is nonnegative,
\[
 \E\left[\left|f_\xi-\frac12\right|\,\middle|\,\mathcal A\right]
 \leq \frac{\E|f_\xi-1/2|}{\P(\mathcal A)}.
\]
This proves the same bound under either conditioning.
\end{proof}

\subsection{Counting centers on one face}

The following projection identity explains the special role of half
rank.

\begin{lemma}[Complementary images]
\label{lem:transverse-projection}
Suppose $B\in\mathcal S_{2r}$ is nonsingular and $A\in\mathcal S_{2r}$
has rank $r$. Then
\begin{equation}
\label{eq:transverse-projection}
 \rank(B-A)=r\quad\Longleftrightarrow\quad AB^{-1}A=A.
\end{equation}
If $A$ and $A-B$ have discriminant signs $\xi_1,\xi_2$, respectively,
under these conditions, then
\begin{equation}
\label{eq:transverse-pair-sign}
 \chi(\det B)=c_0\xi_1\xi_2.
\end{equation}
\end{lemma}

\begin{proof}
If $A$ and $B-A$ both have rank $r$, their images span the image of
$B$, so these images are complementary. Since
$AB^{-1}+(B-A)B^{-1}$ is the identity, uniqueness of the corresponding
direct-sum decomposition makes $AB^{-1}$ the projection onto the first
image. In particular $AB^{-1}A=A$. Conversely, this identity makes
$AB^{-1}$ an idempotent of rank $r$. Its complementary idempotent has
rank $r$ and, after multiplication by $B$, equals $B-A$.

Choose basis-column matrices $P_1,P_2$ of the two complementary images
and write $A=P_1R_1P_1^{\mathsf T}$ and
$A-B=P_2R_2P_2^{\mathsf T}$. Then
\[
 B=(P_1\ P_2)\diag(R_1,-R_2)(P_1\ P_2)^{\mathsf T}.
\]
Taking determinant characters gives
Equation~\eqref{eq:transverse-pair-sign}.
\end{proof}

For independent point matrices $(X_i)_{i\in I}$ define
\begin{equation}
\label{eq:transverse-event}
 \eta_e^X=\chi(\det(X_i-X_k))\quad(e=\{i,k\}),\qquad
 \mathcal T=\{X_i-X_k\text{ nonsingular for every }e\in E\}.
\end{equation}
The choice of orientation of $e$ is immaterial because $D$ is even.
Each difference is uniform in $\mathcal S_D$, so
$\P(\mathcal T^c)=O(q^{-1})$ by Lemma~\ref{lem:symmetric-counts}.
For an activity pattern $B\subseteq\mathscr C$, let
$A_j=\{i\in j:(i,j)\in B\}$ and define the local center factor and its
sign model by
\begin{align}
 F_j(X)&=\E_{Y_j}\prod_{i\in A_j}K_{\xi_{ij}}(X_i-Y_j),
 \label{eq:transverse-face-factor}\\
 G_j(X)&=\prod_{\{i,k\}\subseteq A_j}
       (1+c_0\eta^X_{\{i,k\}}\xi_{ij}\xi_{kj}).
 \label{eq:transverse-model}
\end{align}
On transverse local configurations, $G_j$ equals $2$ for two active
incidences whose pair signs match, and $8$ for three active incidences
whose three pair signs match; it is zero when any required match fails.
With at most one active incidence, both $F_j$ and $G_j$ equal $1$.

\begin{lemma}[Local transverse approximation]
\label{lem:local-transverse}
Suppose $r>1$. For each face $j$, let $\mathcal T_j$ be the event that
all three differences between its point matrices are nonsingular.
Then, uniformly over the activity and sign choices,
\begin{equation}
\label{eq:transverse-local-l1}
 \E\bigl[\1_{\mathcal T_j}|F_j-G_j|\bigr]=O(q^{-1/2}).
\end{equation}
Moreover $0\leq F_j\leq C_D$ on $\mathcal T_j$, for a constant $C_D$
independent of $q$, and $0\leq G_j\leq8$ everywhere.
\end{lemma}

\begin{proof}
First suppose all three incidences are active, and locally write their
point matrices as $X_1,X_2,X_3$ and their signs as $\xi_1,\xi_2,\xi_3$.
Set
\[
 B_0=X_1-X_2,\quad B_1=X_1-X_3,\quad A=X_1-Y,
 \qquad Q=B_0^{-1},\quad L=B_0^{-1}-B_1^{-1}.
\]
On $\mathcal T_j$, the form
$L=B_0^{-1}(B_1-B_0)B_1^{-1}$ is symmetric and nonsingular. Write a
rank-$r$ matrix $A$ as $PRP^{\mathsf T}$. By
Lemma~\ref{lem:transverse-projection}, its other two rank conditions are
equivalent to
\[
 R(P^{\mathsf T}B_\nu^{-1}P)R=R\quad(\nu=0,1),
\]
or to
\begin{equation}
\label{eq:transverse-center-parametrization}
 P^{\mathsf T}LP=0,\qquad
 P^{\mathsf T}QP\text{ nonsingular},\qquad
 R=(P^{\mathsf T}QP)^{-1}.
\end{equation}
Thus the centers satisfying the three rank conditions are in bijection
with the maximal $L$-isotropic subspaces $U$ on which $Q$ restricts
nonsingularly. Changing the basis of $U$ leaves the resulting $A$
unchanged, and every $A$ recovers $U$ as its image, proving the stated
bijection.

Let $\mathcal M$ be the event that all three required pair matches
$\eta^X_{\{a,b\}}=c_0\xi_a\xi_b$ hold. If any match fails, the signed
count is zero by Equation~\eqref{eq:transverse-pair-sign}. If they all
hold, then
\[
 \chi(\det L)
 =\eta^X_{\{1,2\}}\eta^X_{\{1,3\}}\eta^X_{\{2,3\}}
 =c_0,
\]
so $L$ is split by Lemma~\ref{lem:transverse-split}. In
Equation~\eqref{eq:transverse-center-parametrization}, the sign of $A$
is the sign of $Q|_U$, since inversion preserves a nonzero determinant
character. Once it equals $\xi_1$, the pair matches with point $1$
force the other two signs to equal $\xi_2,\xi_3$.

For split $L$, let $f_{\xi_1}(Q;L)$ denote the fraction of its $N_r$
maximal isotropic subspaces having a nonsingular restriction of sign
$\xi_1$. On a matching transverse triple the exact normalized count is
\begin{equation}
\label{eq:transverse-three-normalization}
 F_j=\frac{q^{-b(D)}N_r}{\pi_{\xi_1}\pi_{\xi_2}\pi_{\xi_3}}
       f_{\xi_1}(Q;L)
     =(16+O(q^{-1}))f_{\xi_1}(Q;L).
\end{equation}
Here the exponent cancels because
$r(r-1)/2-b(2r)+3b(r)=0$. The factor $16$ combines the leading
factor $2$ in $N_r$ with the three signed-layer normalization factors
$2$; the restriction-sign fraction will supply a further factor $1/2$.

We specify the probability law used for that fraction. Before any
invertibility conditioning, $B_0,B_1$ are independent uniform symmetric
matrices: this follows by conditioning on $X_1$ and translating the
independent matrices $X_2,X_3$. After requiring $B_0,B_1$ to be
nonsingular, inversion preserves their independent uniform laws on
nonsingular symmetric matrices. Conditional also on a fixed value of
$L=B_0^{-1}-B_1^{-1}$, the matrix $Q$ is therefore uniform on
\[
 \{Q\in\mathcal S_D:Q\text{ and }Q-L\text{ are nonsingular}\}.
\]
Each admissible $Q$ corresponds to exactly one ordered inverse pair
$(Q,Q-L)$, which verifies the conditional uniformity directly.
Requiring $L$ to be nonsingular is exactly the remaining transverse
condition. Lemma~\ref{lem:transverse-restrictions}, uniformly for every
fixed nonsingular split $L$, now gives
\begin{equation}
\label{eq:transverse-matching-error}
 \E\left[\1_{\mathcal M}
       \left|f_{\xi_1}(Q;L)-\frac12\right|
       \,\middle|\,\mathcal T_j,L\right]
 \leq \E\left[\left|f_{\xi_1}(Q;L)-\frac12\right|
       \,\middle|\,\mathcal T_j,L\right]
 =O(q^{-1/2}).
\end{equation}
We used $r>1$, for which $N_r^{-1/2}=O(q^{-1/2})$.
The matching indicator is kept inside the expectation; no additional
conditioning on the matches is used. For nonsplit $L$, the matching
event is empty and both the signed count and its model vanish.
Equations~\eqref{eq:transverse-three-normalization}
and~\eqref{eq:transverse-matching-error} prove the local error estimate
with limiting value $8$ on matches. The same exact formula, with
$0\leq f_{\xi_1}\leq1$, gives a uniform bound on every transverse
triple, including the exceptional restrictions.

For two active incidences, relabel them $1,2$ and put $B=X_1-X_2$ and
$Q=B^{-1}$. The same projection argument now parametrizes centers by
all $r$-subspaces $U$ for which $Q|_U$ is nonsingular. On the pair match
$\chi(\det B)=c_0\xi_1\xi_2$, only the sign $\xi_1$ of this restriction
needs to be imposed; off the match the count is zero. If $f_{\xi_1}(Q)$
is the fraction among all $\genfrac{[}{]}{0pt}{}{2r}{r}_q$ subspaces,
the exact count on the match is
\begin{equation}
\label{eq:transverse-two-normalization}
 F_j=\frac{q^{-b(D)}\genfrac{[}{]}{0pt}{}{2r}{r}_q}
              {\pi_{\xi_1}\pi_{\xi_2}}
       f_{\xi_1}(Q)
     =(4+O(q^{-1}))f_{\xi_1}(Q).
\end{equation}
The exponent identity here is $r^2-b(2r)+2b(r)=0$. Conditional on $B$
being nonsingular, $Q$ is uniform nonsingular symmetric. Apply
Lemma~\ref{lem:transverse-restrictions} with this single conditioning
and keep the matching indicator inside the expectation as before.
This gives limiting value $2$, error $O(q^{-1/2})$, and a uniform
bound on the nonsingular-pair event. Restricting further to
$\mathcal T_j$ can only decrease the nonnegative error expectation.
The remaining cases follow from the mean-one normalization of each
$K_\xi$ and translation invariance.
\end{proof}

\subsection{Joint distribution of the pair signs}

The pairs share point matrices, so marginal balance of their signs does
not give their joint distribution. For each nonempty subset moment,
we first establish a high-probability event for the difference minors
at a selected incident vertex, then condition on all entries except
its off-diagonal row vector and transpose and average over that vector.
Table~\ref{tab:complex} shows that the degree of each
vertex in $(I,E)$ is at most $6$:
its degree there equals its number of incident faces, since each face
contains two pairs through that point and each pair occurs twice.

\begin{lemma}[Joint determinant signs]
\label{lem:joint-signs}
For fixed even $D\geq12$, the law of $(\eta_e^X)_{e\in E}$ converges
in total variation to that of independent uniform signs on
$\{+1,-1\}^E$.
\end{lemma}

\begin{proof}
Fix a nonempty subset $F\subseteq E$, and choose a vertex $i$ incident
with $s\geq1$ of its pairs. Write $N=D-1$ and, for every point $v$,
write its matrix in block form
\[
 X_v=\begin{pmatrix}a_v&z_v^{\mathsf T}\\z_v&W_v\end{pmatrix},
 \qquad W_v\in\mathcal S_N,\quad z_v\in\F_q^N.
\]
Let $k_1,\ldots,k_s$ be the selected neighbors of $i$, and put
$M_\ell=W_i-W_{k_\ell}$. We first establish a high-probability event
using these lower-right blocks, before conditioning on the other
entries. Given $W_i$, the $M_\ell$ are independent uniform symmetric
matrices. All are nonsingular except with probability $O(q^{-1})$.
Conditional on their nonsingularity, their inverses are independent
uniform nonsingular symmetric matrices.

Let $Q_1,\ldots,Q_s$ instead be independent uniform matrices in
$\mathcal S_N$. For any $t\in\F_q^s\setminus\{0\}$, the matrix
$\sum_\ell t_\ell Q_\ell$ is uniform in $\mathcal S_N$: condition on
all but a coordinate with nonzero coefficient. By
Equation~\eqref{eq:transverse-nullity}, its probability of rank at most
$s$ is $O(q^{-b(N-s)})$. The union bound over all nonzero $t$ gives
\[
 \P\left(\exists t\ne0:
       \rank\left(\sum_\ell t_\ell Q_\ell\right)\leq s\right)
 =O(q^{s-b(N-s)}).
\]
Conditioning all the $Q_\ell$ to be nonsingular divides this upper bound
by a probability $1-O(q^{-1})$. It follows that the event
\begin{equation}
\label{eq:transverse-good-minors}
 \mathcal G=\left\{\begin{array}{l}
 M_\ell\text{ is nonsingular for every }\ell,\\
 \rank(\sum_\ell t_\ell M_\ell^{-1})\geq s+1
       \text{ for every }t\in\F_q^s\setminus\{0\}
 \end{array}\right\}
\end{equation}
satisfies
\begin{equation}
\label{eq:transverse-good-probability}
 \P(\mathcal G^c)=O(q^{-1}+q^{s-b(N-s)})=O(q^{-1}).
\end{equation}
Indeed $s\leq6$ and $D\geq12$ imply $N-s\geq5$ and
$b(N-s)\geq15>s$. These probability bounds hold given every $W_i$
and hence hold without conditioning as well.

Only now condition on all entries of all point matrices except $z_i$.
The event $\mathcal G$ is determined by this conditioning, whereas
$z_i$ remains uniform on $\F_q^N$. On every such fixed configuration
in $\mathcal G$, the determinant formula gives
\begin{equation}
\label{eq:transverse-schur-evaluations}
 \det(X_i-X_{k_\ell})=\det(M_\ell)h_\ell(z_i),\qquad
 h_\ell(z)=a_i-a_{k_\ell}
 -(z-z_{k_\ell})^{\mathsf T}M_\ell^{-1}(z-z_{k_\ell}).
\end{equation}
Let $\mu$ be the law of $(h_1(z),\ldots,h_s(z))$ for uniform $z$.
Fix a nontrivial additive character $\psi$ of $\F_q$ and define
$\widehat\mu(t)=\E_z\psi(\sum_\ell t_\ell h_\ell(z))$.
The quadratic part at every frequency $t\ne0$ has rank at least $s+1$
by Equation~\eqref{eq:transverse-good-minors}.

We use the standard rank bound for a quadratic Gauss sum; see
Gowers and Wolf~\cite[Lemma~3.3]{gowerswolf2011}. The short differencing argument
also fixes our normalization. A quadratic polynomial
$p(z)=z^{\mathsf T}Hz+\lambda^{\mathsf T}z+c$ with symmetric quadratic
matrix $H$ obeys
\begin{align*}
 \left|\E_z\psi(p(z))\right|^2
 &=\E_{z,u}\psi(p(z+u)-p(z))\\
 &=\E_u\psi(u^{\mathsf T}Hu+\lambda^{\mathsf T}u)
               \E_z\psi(2z^{\mathsf T}Hu),
\end{align*}
whose absolute value is at most
$\P_u(Hu=0)=q^{-\rank H}$. Here the inner average vanishes unless
$Hu=0$, by additive-character orthogonality and the fact that $q$ is
odd. Therefore
$|\widehat\mu(t)|\leq q^{-(s+1)/2}$ for $t\ne0$, uniformly over all
the fixed linear and constant terms in
Equation~\eqref{eq:transverse-schur-evaluations}.

Writing $\nu$ for uniform measure on $\F_q^s$, finite Fourier
orthogonality and Cauchy--Schwarz yield the explicit bound
\begin{align}
\label{eq:transverse-fourier-tv}
 \TV(\mu,\nu)
 &=\frac12\sum_x|\mu(x)-q^{-s}|\\
 &\leq\frac12\left(q^s\sum_x|\mu(x)-q^{-s}|^2\right)^{1/2}
 =\frac12\left(\sum_{t\ne0}|\widehat\mu(t)|^2\right)^{1/2}
 \leq\frac12q^{-1/2}.\nonumber
\end{align}
The identity in the middle follows by expanding the square and using
$\sum_t\psi(t\cdot(x-y))=q^s\1_{\{x=y\}}$.

Under the same fixed conditioning, all factors
$\eta_e^X$ for $e\in F$ not incident to $i$ are constant. They and the
characters of the minor determinants have absolute value at most one.
By Equation~\eqref{eq:transverse-schur-evaluations}, the remaining part
of the product is $\prod_{\ell=1}^s\chi(h_\ell(z_i))$. Its expectation
under $\nu$ is zero, because $\E_{x\in\F_q}\chi(x)=0$ and $s\geq1$.
The difference of expectations of a function bounded by one is at
most twice total variation, so on $\mathcal G$ the absolute conditional
expectation of the entire selected product is at most $q^{-1/2}$.
On $\mathcal G^c$ it is at most one. Averaging and using
Equation~\eqref{eq:transverse-good-probability} proves
\begin{equation}
\label{eq:transverse-subset-moments}
 \E\prod_{e\in F}\eta_e^X=O(q^{-1/2})\qquad(F\ne\varnothing).
\end{equation}

Finally, replace any zero coordinate of $\eta^X$ by $+1$, obtaining a
sign vector $\widetilde\eta^X$. A union bound and
Lemma~\ref{lem:symmetric-counts} show that this replacement changes
the vector with probability $O(q^{-1})$ and changes each subset moment
by $O(q^{-1})$. For every $\varepsilon\in\{+1,-1\}^E$, the identity
\[
 \1_{\{\widetilde\eta^X=\varepsilon\}}
 =2^{-|E|}\prod_{e\in E}(1+\varepsilon_e\widetilde\eta_e^X)
\]
expresses its probability through the subset moments. The empty
subset contributes $2^{-|E|}$ and every other subset is $O(q^{-1/2})$
by Equation~\eqref{eq:transverse-subset-moments}. Since $E$ is fixed,
summing over all sign vectors proves convergence in total variation.
\end{proof}

\subsection{Assembling the transverse contribution}

\begin{proposition}
\label{prop:transverse-limit}
For every fixed even $D\geq12$, every activity pattern
$B\subseteq\mathscr C$, and every prescribed set of incidence signs,
\begin{equation}
\label{eq:transverse-limit}
 \E\left[\1_{\mathcal T}
       \prod_{(i,j)\in B}K_{\xi_{ij}}(X_i-Y_j)\right]
 =\E_\eta\prod_{j\in J}\prod_{\{i,k\}\subseteq A_j}
       (1+c_0\eta_{\{i,k\}}\xi_{ij}\xi_{kj})+o(1),
\end{equation}
where the $\eta_e$ are independent uniform signs. At each fixed $D$,
the error is uniform over all activity patterns and incidence signs.
\end{proposition}

\begin{proof}
Given the point matrices, the center variables are independent, so the
left side is $\E[\1_{\mathcal T}\prod_jF_j]$. Increase $C_D$ from
Lemma~\ref{lem:local-transverse} to at least $8$. The deterministic
telescoping identity for two products gives, on $\mathcal T$,
\[
 \left|\prod_jF_j-\prod_jG_j\right|
 \leq C_D^{|J|-1}\sum_j|F_j-G_j|.
\]
Since $\mathcal T\subseteq\mathcal T_j$, taking expectations and
applying Equation~\eqref{eq:transverse-local-l1} bounds this difference
by $O(q^{-1/2})$. This argument applies to faces sharing point matrices;
it uses no independence of their local errors.

The bound $\prod_jG_j\leq8^{|J|}$ and
$\P(\mathcal T^c)=O(q^{-1})$ allow removal of
$\1_{\mathcal T}$ from the model expectation. Lemma~\ref{lem:joint-signs}
then replaces $\eta^X$ by independent uniform signs. Total variation
controls every bounded function, including the displayed model even
when $c_0$ varies with the prime $q$. Finally, there are only finitely
many activity and sign choices, so the largest of their errors also
tends to zero.
\end{proof}

Proposition~\ref{prop:transverse-limit} reduces
Lemma~\ref{lem:rank-layer} to the weighted singular contribution
\begin{equation}
\label{eq:transverse-tail-target}
 \E\left[\1_{\mathcal T^c}
       \prod_{(i,j)\in B}K_{\xi_{ij}}(X_i-Y_j)\right]=o(1).
\end{equation}
The unweighted estimate $\P(\mathcal T^c)=O(q^{-1})$ does not establish
Equation~\eqref{eq:transverse-tail-target}, because the normalized
rank-layer indicators grow with $q$. Proposition~\ref{prop:singular-tail}
will prove this estimate for all sufficiently large even $D$, after
the geometric and exponent estimates in the next sections. The
threshold $D\geq12$ above applies only to the transverse calculation
and the joint-sign argument.
